% Figure 2.14 (fig:gauss). Reference triangle: the nodes carry the
% displacements and the nodal forces, the three Gauss points the strains and
% the stresses (notation of the discrete Tonti diagram, fig:tonti-fem).
\begin{tikzpicture}[scale=4,>={Latex[length=2mm]},
  nd/.style={circle,fill=black,inner sep=1.7pt},
  gp/.style={circle,fill=red!75!black,inner sep=1.7pt}]
  % axes and the reference element
  \draw[->] (-0.1,0) -- (1.3,0) node[below] {$\xi_1$};
  \draw[->] (0,-0.1) -- (0,1.25) node[left] {$\xi_2$};
  \fill[cu!15] (0,0) -- (1,0) -- (0,1) -- cycle;
  \draw[thick] (0,0) -- (1,0) -- (0,1) -- cycle;
  \node[below left] at (0,0) {$0$};
  \node[below] at (1,0) {$1$};
  \node[left] at (0,1) {$1$};
  \foreach \p in {(0,0),(1,0),(0,1)} {\node[nd] at \p {};}
  \foreach \p in {(1/6,1/6),(2/3,1/6),(1/6,2/3)} {\node[gp] at \p {};}
  % the two operators
  \draw[->,cu,thick,shorten <=3pt,shorten >=3pt] (1,0) to[out=80,in=40,looseness=1.6]
    node[pos=0.5,above right=-1pt,font=\small] {$\mat{B}$} (2/3,1/6);
  \draw[->,cphi!85!black,thick,shorten <=3pt,shorten >=3pt] (1/6,2/3) to[out=50,in=10,looseness=1.6]
    node[pos=0.5,above right=-1pt,font=\small] {$\mat{B}\T$} (0,1);
  % legend
  \begin{scope}[shift={(1.35,0.95)},font=\small]
    \node[nd] at (0,0) {};
    \node[right=4pt,align=left] at (0,0)
      {node: \ \textcolor{cu}{$\mat{u}$}, \ \textcolor{cphi}{$\mat{F}$}};
    \node[gp] at (0,-0.17) {};
    \node[right=4pt,align=left] at (0,-0.17)
      {Gauss point: \ \textcolor{cu}{$\mat{e}$}, \ \textcolor{cphi}{$\mat{s}$}};
    \node[right,cu] at (-0.03,-0.42) {$\mat{e}=\mat{B}\mat{u}$ \ (nodes $\to$ Gauss points)};
    \node[right,cphi!85!black] at (-0.03,-0.6) {$\mat{F}=\mat{B}\T\mat{s}$ \ (Gauss points $\to$ nodes)};
  \end{scope}
\end{tikzpicture}
